\section{Mapping Research Questions, Objectives, and Contributions}
\label{sec:rq_ro_mapping}

Table \ref{tab:rq_ro_mapping} illustrates the alignment between the research questions, specific objectives, the evaluation metrics collected, and the final contributions of the project. This structured mapping ensures that the evaluation effectively addresses the core problem statement.

\begin{table}[H]
\centering
\caption{Mapping of Research Questions to Objectives, Metrics, and Contributions}
\label{tab:rq_ro_mapping}
\small
\begin{tabularx}{\textwidth}{>{\hsize=0.6\hsize}X >{\hsize=0.8\hsize}X >{\hsize=1.2\hsize}X >{\hsize=1.4\hsize}X}
\toprule
\textbf{Research Question} & \textbf{Related Objective} & \textbf{Evaluation Focus} & \textbf{Related Contribution} \\
\midrule
\textbf{RQ1:} How can a modular research-grade benchmarking framework be designed? & \textbf{RO1:} Design and implement a modular research-grade ZK-rollup benchmarking framework. & Architectural separation, integration of microservices, local testnet deployment success. & \textbf{C1:} A modular ZK-rollup benchmarking framework with separable components. \\
\addlinespace
\textbf{RQ2:} How do configurations affect throughput, latency, and cost? & \textbf{RO2:} Develop a reproducible experimental pipeline. \newline \textbf{RO3:} Empirically evaluate performance trade-offs. & Throughput, latency, finalization time, L1 gas cost, DA footprint, proving cycles under configuration sweeps. & \textbf{C2:} A reproducible experimental pipeline. \newline \textbf{C3:} Empirical findings on throughput-latency-cost trade-offs. \\
\addlinespace
\textbf{RQ3:} What practical implementation bottlenecks emerge? & \textbf{RO3:} Empirically evaluate performance trade-offs (specifically isolating hardware and queueing bottlenecks). & Proof generation overhead (wall-clock, cycles), sequencer fairness/ordering anomalies, telemetry gaps. & \textbf{C3:} Empirical findings including anomalies such as RISC0 step-function behavior and nonce ordering issues. \\
\bottomrule
\end{tabularx}
\end{table}
